\documentclass[11pt]{article}
\usepackage[T1]{fontenc}
\usepackage[utf8]{inputenc}
\usepackage[french]{babel}
\usepackage{lmodern}
\usepackage{microtype}
\usepackage[margin=2.3cm]{geometry}
\usepackage{booktabs,tabularx,longtable,array}
\usepackage{enumitem}
\usepackage[dvipsnames]{xcolor}
\usepackage[most]{tcolorbox}
\usepackage{hyperref}
\hypersetup{colorlinks=true,linkcolor=MidnightBlue,urlcolor=MidnightBlue}

\setlist{itemsep=2pt,topsep=4pt}
\newcolumntype{L}[1]{>{\raggedright\arraybackslash}p{#1}}
\newcolumntype{Y}{>{\raggedright\arraybackslash}X}

\newtcolorbox{encadre}[1]{colback=MidnightBlue!4,colframe=MidnightBlue!70,
  fonttitle=\bfseries,title={#1},arc=1mm,boxrule=0.6pt,left=6pt,right=6pt}

\title{\textbf{STT-7360 Méthodes biostatistiques}\\[2pt]
\large Proposition de plan pour la session d'hiver 2027\\
\normalsize Cours orienté vers la formation à la recherche (maîtrise et doctorat)}
\author{Julien Miron}
\date{Version de travail --- \today}

\begin{document}
\maketitle

\begin{encadre}{En bref}
Le plan conserve l'ossature du cours d'hiver 2026 (rappels, simulations, données corrélées,
théorie semi-paramétrique, données manquantes, erreur de mesure et échantillonnage à deux phases,
tests multiples, trois études de cas, présentations d'articles). Il déplace toutefois le centre de
gravité de l'\emph{acquisition de méthodes} vers l'\emph{exercice du métier de chercheur} :
lire un article méthodologique de façon critique, en reproduire les résultats, évaluer une méthode
par simulation, rédiger, arbitrer et présenter. L'examen final est remplacé (ou allégé) par un
mini-projet de recherche qui s'étend sur toute la session.
\end{encadre}

\section{Objectifs de formation}

À la fin du cours, l'étudiante ou l'étudiant devrait pouvoir :
\begin{enumerate}[label=\textbf{O\arabic*.}]
  \item \textbf{Lire} un article méthodologique en biostatistique et en extraire la structure
        (problème, modèle, estimateur, hypothèses, propriétés, validation) à l'aide d'une
        \emph{fiche des caractéristiques méthodologiques}.
  \item \textbf{Reproduire} les résultats numériques d'un article (analyse de données et
        simulations) dans un flux de travail reproductible (R, Quarto/R Markdown, Git).
  \item \textbf{Concevoir} une étude de simulation rigoureuse (cadre ADEMP : objectifs,
        mécanisme générateur, estimand, méthodes, mesures de performance) et en rapporter
        l'erreur Monte Carlo.
  \item \textbf{Mobiliser} les outils théoriques transversaux (fonction d'influence, équations
        d'estimation, estimateurs sandwich, efficacité semi-paramétrique) pour comprendre
        \emph{pourquoi} une méthode fonctionne et quand elle échoue.
  \item \textbf{Choisir et justifier} une méthode pour des données corrélées, incomplètes,
        mesurées avec erreur ou soumises à de la multiplicité.
  \item \textbf{Communiquer} : rédiger un court article, produire un rapport d'arbitrage et
        présenter oralement un travail de recherche.
\end{enumerate}

\section{Principes pédagogiques}

\begin{itemize}
  \item \textbf{Un fil conducteur : « de l'article à la méthode, de la méthode à la question
        ouverte ».} Chaque bloc thématique se termine par une question de recherche plausible
        (limite d'une méthode, hypothèse non vérifiée, extension naturelle).
  \item \textbf{La simulation comme outil de recherche}, introduite dès la 2\textsuperscript{e}
        semaine et réutilisée dans toutes les études de cas et dans le projet.
  \item \textbf{La théorie comme langage commun.} La fonction d'influence sert de pont entre la
        robustesse, les équations d'estimation généralisées (GEE), la pondération par
        l'inverse de la probabilité et les méthodes ciblées (TMLE).
  \item \textbf{Deux parcours} (comme en 2026) : un tronc commun, puis des lectures et exercices
        supplémentaires plus théoriques pour le doctorat (Tsiatis, van der Laan et Rose,
        Jiang et Nguyen).
  \item \textbf{Classe active} : lectures préalables, courte présentation magistrale,
        puis discussion d'article ou atelier de programmation.
\end{itemize}

\section{Calendrier proposé (15 semaines)}

Les semaines sont numérotées ; la semaine de lecture est placée en semaine 9.
\emph{M} = tronc commun (maîtrise et doctorat), \emph{D} = complément doctorat.

\renewcommand{\arraystretch}{1.25}
\begin{longtable}{@{}c L{3.6cm} L{6.3cm} L{4.2cm}@{}}
\toprule
\textbf{Sem.} & \textbf{Thème} & \textbf{Contenu et activités} & \textbf{Volet recherche / livrable}\\
\midrule
\endhead
\bottomrule
\endfoot
1 & Introduction et rappels &
Vraisemblance, modèles linéaires généralisés, inférence asymptotique (notes \emph{Rappel}).
Présentation du projet de session. &
Atelier : flux de travail reproductible (Git, Quarto, \texttt{renv}).\\
2 & Simulations &
Principes, générateurs, cadre ADEMP, mesures de performance (biais, EQM, couverture),
erreur Monte Carlo, parallélisation. &
Exercice : critiquer la section « simulations » d'un article publié.\\
3 & Données corrélées complexes (I) &
Modèles mixtes linéaires et généralisés ; effet conditionnel vs marginal (Wu, ch.~2 ;
Fitzmaurice et coll.). \emph{D} : GAM mixtes, Jiang et Nguyen 3.5.2. &
Devoir 1 (modèles mixtes, données BSI).\\
4 & Données corrélées complexes (II) &
GEE, estimateur sandwich, comparaison GEE/GLMM (articles Zhang et coll.). &
Choix des articles à présenter.\\
5 & \textbf{Étude de cas 1} &
Neuhaus et McCulloch (2006) : corrélation entre effets aléatoires et covariables ;
régression logistique conditionnelle (données \texttt{naisga}). &
Fiche méthodologique + reproduction d'un tableau de l'article.\\
6 & Théorie semi-paramétrique (I) &
Fonction d'influence, M-estimateurs, robustesse (Wu, ch.~9 ; Tsiatis, ch.~3).
Application : modèle mixte robuste (\texttt{SleepStudyRobuste.R}). &
\textbf{Proposition de projet} (2 pages).\\
7 & Théorie semi-paramétrique (II) &
Équations d'estimation, GEE vues comme M-estimation, efficacité.
\emph{D} : Tsiatis, ch.~4 ; van der Laan et Rose, ch.~5. &
Devoir 2 (GEE, pleuropneumonie).\\
8 & \textbf{Étude de cas 2} &
Méthodes ciblées pour variables continues (Tuglus et van der Laan) ; TMLE, découverte de
biomarqueurs. &
Fiche méthodologique + petite simulation comparant TMLE et régression.\\
9 & \multicolumn{3}{l}{\emph{Semaine de lecture} --- rencontres individuelles sur les projets}\\
10 & Données manquantes (I) &
Mécanismes MCAR/MAR/MNAR, imputation multiple (Wu, sect.~3.2--3.5). &
Exercice d'imputation (données dentaires).\\
11 & Données manquantes (II) &
Pondération par l'inverse de la probabilité (IPW), estimateurs doublement robustes ;
lien avec la fonction d'influence (Wu 10.3.1). &
\textbf{Rapport d'étape} du projet (résultats préliminaires).\\
12 & \textbf{Étude de cas 3} &
Erreur de mesure et échantillonnage à deux phases (Breslow et Chatterjee, 1999 ;
données NWTS). Comparaison d'estimateurs. &
Fiche méthodologique + rapport d'arbitrage simulé de l'article.\\
13 & Tests multiples &
FWER, FDR, procédures pas-à-pas, rééchantillonnage (Farcomeni ; Dudoit et van der Laan,
données ApoAI). &
Résolution de problèmes en équipe.\\
14 & Présentations d'articles &
Présentations de 20 min + 10 min de discussion animée par un « arbitre » désigné. &
Rapport d'arbitrage écrit (1--2 pages).\\
15 & Colloque de fin de session &
Présentations des projets (format congrès : 12 min + questions). &
\textbf{Article final} (remis une semaine plus tard).\\
\end{longtable}

\section{Le mini-projet de recherche (fil rouge de la session)}

\begin{encadre}{Format : « reproduire, puis étendre »}
Chaque personne (ou équipe de deux à la maîtrise) choisit un article méthodologique récent,
idéalement lié à son propre sujet de recherche, en reproduit un résultat clé, puis propose
\emph{une} extension évaluée par simulation ou par une petite dérivation théorique.
\end{encadre}

\paragraph{Exemples d'extensions.}
\begin{itemize}
  \item Évaluer la sensibilité d'un estimateur à une violation d'hypothèse (mauvaise
        spécification des effets aléatoires, MNAR, contamination par des valeurs aberrantes).
  \item Comparer la méthode de l'article à une alternative robuste ou doublement robuste.
  \item Dériver la fonction d'influence d'un estimateur et en tirer une variance sandwich,
        puis vérifier la couverture par simulation (\emph{D}).
  \item Adapter la méthode à un autre type de réponse (binaire, ordinale, de comptage).
\end{itemize}

\paragraph{Pistes tirées de la liste d'articles 2026.}
Estimateur sandwich empirique pour la g-computation par espérances conditionnelles itérées
(Talbot et coll.) ; ACP fonctionnelle comme alternative aux modèles mixtes (Ségalas et coll., 2024) ;
modèles longitudinaux binomiaux et de comptage, GEE vs GLMM (Zhang et coll.) ;
modèles mixtes avec erreur de mesure et M-estimateurs pour modèles mixtes (Wu, sect.~5.1--5.2 et 9.4).

\paragraph{Jalons.}
\begin{center}
\begin{tabular}{@{}lll@{}}
\toprule
Semaine & Livrable & Contenu attendu\\
\midrule
6  & Proposition & question, article de départ, estimand, plan de simulation ADEMP\\
9  & Rencontre   & validation du plan et du code de reproduction\\
11 & Rapport d'étape & reproduction réussie ou écarts expliqués\\
15 & Présentation & format congrès\\
16 & Article     & 8--12 pages (maîtrise) ou 12--20 pages (doctorat), code public\\
\bottomrule
\end{tabular}
\end{center}

\section{Études de cas : format révisé}

Les trois études de cas de 2026 sont conservées, mais chacune produit trois éléments courts :
\begin{enumerate}
  \item la \textbf{fiche des caractéristiques méthodologiques} (déjà utilisée en 2026) ;
  \item une \textbf{reproduction} d'un tableau ou d'une figure de l'article à partir des données
        fournies (\texttt{naisga.csv}, \texttt{genes\_select}, données NWTS) ;
  \item une \textbf{question ouverte} : une limite de l'article et une façon de l'étudier.
\end{enumerate}
L'étude de cas 3 prend la forme d'un rapport d'arbitrage, afin d'exercer cette compétence avant
les présentations d'articles.

\section{Évaluation proposée}

\begin{center}
\begin{tabularx}{\linewidth}{@{}Y r Y@{}}
\toprule
\textbf{Élément} & \textbf{Poids} & \textbf{Remarques}\\
\midrule
Devoirs (2) & 20\,\% & modèles mixtes ; GEE et sandwich\\
Études de cas (3) & 25\,\% & fiche + reproduction + question ouverte\\
Présentation d'article et rapport d'arbitrage & 15\,\% & évaluation croisée en partie\\
Projet de recherche & 40\,\% & proposition 5\,\%, rapport d'étape 5\,\%, présentation 10\,\%, article 20\,\%\\
\bottomrule
\end{tabularx}
\end{center}

\noindent\emph{Variante :} si l'on souhaite garder un examen final (comme en 2026), il peut
remplacer l'article final (20\,\%) et porter surtout sur la théorie transversale (fonction
d'influence, équations d'estimation, mécanismes de données manquantes). Les questions de pratique
et l'examen 2023 existants peuvent servir de base.

\paragraph{Différenciation maîtrise / doctorat.}
Au doctorat : exercices théoriques supplémentaires (déjà prévus en 2026), article de projet plus
long avec une composante théorique, et rôle d'arbitre principal lors des présentations.

\section{Compétences de recherche transversales}

\begin{center}
\begin{tabularx}{\linewidth}{@{}L{4.2cm} Y@{}}
\toprule
\textbf{Compétence} & \textbf{Où elle est travaillée}\\
\midrule
Lecture critique & fiches méthodologiques, études de cas, présentations\\
Reproductibilité & atelier de la semaine 1, reproductions, dépôt de code du projet\\
Simulation & semaine 2, études de cas 2 et 3, projet\\
Théorie & semaines 6--7 et 11 ; exercices doctorat\\
Rédaction scientifique & proposition, rapport d'étape, article final\\
Arbitrage & étude de cas 3, présentations d'articles\\
Communication orale & présentations d'articles, colloque de fin de session\\
\bottomrule
\end{tabularx}
\end{center}

\section{Matériel à réutiliser ou à préparer}

\begin{itemize}
  \item \textbf{Réutilisable tel quel} (site Brio 2026) : notes \emph{Rappel}, \emph{Simulations},
        \emph{Équations d'estimation}, \emph{Données manquantes}, exercices et solutionnaires
        (\texttt{.tex}, \texttt{.Rmd}), données des études de cas, examen 2023 et questions de pratique.
  \item \textbf{À adapter} : notes de cours en \texttt{.docx} (cours 3, semi-paramétrique,
        données manquantes, cours 10) à convertir en LaTeX/Quarto pour uniformiser ;
        mise à jour de la liste d'articles pour présentation.
  \item \textbf{À créer} : guide du projet (gabarit d'article et grille), gabarit de rapport
        d'arbitrage, gabarit ADEMP, grilles d'évaluation des présentations, gabarit de dépôt Git.
\end{itemize}

\section*{Références principales}
\begin{itemize}[leftmargin=*,label={}]\small
  \item Dudoit, S. et van der Laan, M. J. (2008). \emph{Multiple Testing Procedures with Applications to Genomics}. Springer.
  \item Fitzmaurice, G. M., Laird, N. M. et Ware, J. H. (2011). \emph{Applied Longitudinal Analysis}, 2\textsuperscript{e} éd. Wiley.
  \item Jiang, J. et Nguyen, T. (2021). \emph{Linear and Generalized Linear Mixed Models and Their Applications}, 2\textsuperscript{e} éd. Springer.
  \item Morris, T. P., White, I. R. et Crowther, M. J. (2019). Using simulation studies to evaluate statistical methods. \emph{Statistics in Medicine}, 38, 2074--2102.
  \item Tsiatis, A. A. (2006). \emph{Semiparametric Theory and Missing Data}. Springer.
  \item van der Laan, M. J. et Rose, S. (2011). \emph{Targeted Learning}. Springer.
  \item Wu, L. (2010). \emph{Mixed Effects Models for Complex Data}. Chapman \& Hall/CRC.
  \item Articles des études de cas : Neuhaus et McCulloch (2006, \emph{JRSS B}) ; Tuglus et van der Laan (2011) ; Breslow et Chatterjee (1999, \emph{JRSS C}).
\end{itemize}

\end{document}
